\documentclass[12pt,a4paper,UTF8]{ctexart}
\usepackage[a4paper,top=22mm,bottom=22mm,left=22mm,right=22mm,headheight=15pt]{geometry}
\usepackage{amsmath,amssymb,mathtools,bm}
\usepackage{graphicx}
\usepackage{xcolor}
\usepackage{booktabs,longtable,array}
\usepackage{enumitem}
\usepackage{fancyhdr}
\usepackage{caption}
\usepackage{float}
\usepackage{microtype}
\usepackage{setspace}
\usepackage[hidelinks,unicode]{hyperref}
\usepackage[most]{tcolorbox}

\definecolor{ContestBlue}{HTML}{17365D}
\definecolor{SoftBlue}{HTML}{EEF4FA}
\definecolor{SoftGray}{HTML}{F5F6F7}
\hypersetup{pdftitle={第100套},pdfauthor={物理竞赛教学资料}}

\setstretch{1.24}
\setlength{\parindent}{2em}
\setlength{\parskip}{0.35em}
\setlist{itemsep=0.35em,topsep=0.4em}
\captionsetup{font=small,labelfont=bf,justification=centering}
\setcounter{secnumdepth}{0}

\pagestyle{fancy}
\fancyhf{}
\fancyhead[L]{\small 第100套}
\fancyfoot[C]{\thepage}
\renewcommand{\headrulewidth}{0.4pt}

\ctexset{
  section={format=\Large\bfseries\color{ContestBlue},beforeskip=1.2em,afterskip=0.7em},
  subsection={format=\large\bfseries\color{ContestBlue},beforeskip=1.0em,afterskip=0.5em},
  subsubsection={format=\normalsize\bfseries,beforeskip=0.8em,afterskip=0.35em}
}

\makeatletter
\newcommand{\tightlist}{\setlength{\itemsep}{0pt}\setlength{\parskip}{0pt}}
\newcommand{\maxwidth}{\ifdim\Gin@nat@width>\linewidth\linewidth\else\Gin@nat@width\fi}
\newcommand{\maxheight}{\ifdim\Gin@nat@height>0.72\textheight 0.72\textheight\else\Gin@nat@height\fi}
\newcommand{\pandocbounded}[1]{\begingroup\setkeys{Gin}{width=\maxwidth,height=\maxheight,keepaspectratio}#1\endgroup}
\makeatother

\begin{document}

\begin{titlepage}
  \centering
  \vspace*{58mm}
  {\Huge\bfseries\color{ContestBlue} 第100套\par}
  \vfill
  \vspace*{28mm}
\end{titlepage}


\section{第 1 题}\label{ux7b2c-1-ux9898}

半径为 \(R\)
的半球形山坡表面光滑，坡顶有一小洞。把高尔夫球视为质点，球最初位于半球外表面上，球心位置对应的半径与竖直向上的半径夹角为
\(\theta\)。现沿球面切线、指向坡顶的方向给球初速度 \(v\)，重力加速度为
\(g\)。

\begin{enumerate}
\def\labelenumi{\arabic{enumi}.}
\tightlist
\item
  若球始终贴着半球表面上滑，并在到达坡顶时速度恰好减为零，求初速度
  \(v\)。
\item
  求第（1）问方案可行时 \(\theta\) 的范围。
\item
  当第（1）问方案不可行时，应怎样使球进入坡顶小洞？求新方案所需的初速度
  \(v\)。
\item
  求第（3）问方案可行时 \(\theta\) 的范围。
\end{enumerate}

\begin{figure}[H]
\centering
\includegraphics[width=0.698\linewidth,height=\textheight,keepaspectratio,alt={图 1-1　小球从半球外表面沿切线方向射向坡顶。}]{../figures/question-01.pdf}
\caption*{图 1-1　小球从半球外表面沿切线方向射向坡顶。}
\end{figure}

\newpage

\section{第 2 题}\label{ux7b2c-2-ux9898}

在坐标原点以初速度 \(v_0\) 斜向上抛出一质点，初速度与水平方向夹角为
\(\theta\)，忽略空气阻力，重力加速度为 \(g\)。

\begin{enumerate}
\def\labelenumi{\arabic{enumi}.}
\tightlist
\item
  质点需要到达点 \((l,0)\)。两点之间有一座山，山体上边界为
  \[y_h(x)=y_0-\frac{k}{2}\left(x-\frac l2\right)^2,\] 且
  \(y_0-kl^2/8<0\)。要求抛物线不与山体相交，求最小初速度，并给出参数之间相应的分类条件。
\item
  另一座山的上边界为 \[y_h(x)=y_0-\frac{1}{2l}(x-x_0)^2,\] 且
  \(y_h(0)<0\)、\(y_h(l)<0\)。要求向右抛出的质点在飞行过程中不与山体相交。求最小初速度及相应的分类条件；以
  \(x_0\) 为横轴、\(y_0\) 为纵轴，标出不同情形所在的参数区域。
\end{enumerate}

\newpage

\section{第 3 题}\label{ux7b2c-3-ux9898}

一单摆由长度为 \(l\)、不带电且不可伸长的细绳和质量为 \(m\)、电量为 \(q\)
的带电质点组成。摆的下方有一接地金属平板。真空静电力常量为
\(k\)，重力加速度为 \(g\)。初始时摆线与竖直方向夹角为
\(\alpha_0=\pi/3\)，质点到金属板的垂直距离为 \(d=3l/4\)。

\begin{enumerate}
\def\labelenumi{\arabic{enumi}.}
\tightlist
\item
  求初始时质点所受合外力的大小和方向。
\item
  求初始时绳子的张力，以及质点第一次到达最低点时绳子的张力。
\item
  质点最终在最低点附近做小振动，求小振动周期 \(f_3\)。
\item
  若把金属平板换成半径为 \(R\) 的接地金属球，其中
  \(R>l\)，且金属球顶点位于悬点在原金属板上的垂足处，求此时最低点附近的小振动周期
  \(f_4\)。
\end{enumerate}

\begin{figure}[H]
\centering
\includegraphics[width=0.761\linewidth,height=\textheight,keepaspectratio,alt={图 3-1　带电单摆位于接地金属板上方。}]{../figures/question-03.pdf}
\caption*{图 3-1　带电单摆位于接地金属板上方。}
\end{figure}

\newpage

\section{第 4 题}\label{ux7b2c-4-ux9898}

用单位长度电阻为 \(r\) 的均匀电阻丝围成图示矩形网络，各线段长度由
\(l\)、\(2l\)、\(4l\) 组成。灰色的 \(4l imes4l\)
中央区域内存在垂直纸面向里的均匀磁场

\[B=kt.\]

\begin{enumerate}
\def\labelenumi{\arabic{enumi}.}
\tightlist
\item
  求支路 \(AB\)、\(CD\) 中的电流。
\item
  按图中虚线接入理想电压表，求电压表读数；读数定义为图中``\(+\)''端电势减去``\(-\)''端电势。
\end{enumerate}

\begin{figure}[H]
\centering
\includegraphics[width=0.593\linewidth,height=\textheight,keepaspectratio,alt={图 4-1　时变磁场中的矩形电阻网络及电压表接法。}]{../figures/question-04.pdf}
\caption*{图 4-1　时变磁场中的矩形电阻网络及电压表接法。}
\end{figure}

\newpage

\section{第 5 题}\label{ux7b2c-5-ux9898}

水平光滑桌面上有两条平行光滑导轨，间距为
\(l\)。导轨所在区域存在垂直桌面向里的匀强磁场 \(B\)，左端用理想电感
\(L\) 连接。质量为 \(m\)
的导体棒横跨两导轨，并与一根``反向弹簧''相连：当棒相对 \(x=0\)
位置的位移为 \(x\) 时，弹簧给棒的力为 \(F=kx\)，方向指向远离 \(x=0\)
的一侧。忽略电阻以及弹簧的一切电磁效应。初始时 \(x=0\)，导体棒以速度
\(v_0\) 向右运动。

\begin{enumerate}
\def\labelenumi{\arabic{enumi}.}
\tightlist
\item
  初始电流为零。若反向弹簧始终保持理想特性，求使导体棒做匀速运动的劲度系数
  \(k\)。
\item
  若 \(k\) 取第（1）问结果的一半，求导体棒的振动周期和平衡位置。
\item
  在第（2）问条件下，求导体棒位移的最大值 \(x_{3\max}\)、最小值
  \(x_{3\min}\)，以及电流绝对值的最大值 \(I_{3m}\)。
\item
  仍取第（1）问劲度系数的一半，但初始时电路中已有大小为 \(I_0\)
  的逆时针电流。把 \(I_0\) 视为小量，求位移最大值 \(x_{4\max}\) 和最小值
  \(x_{4\min}\)，并给出一阶小量结果。
\end{enumerate}

\begin{figure}[H]
\centering
\includegraphics[width=0.9\linewidth,height=\textheight,keepaspectratio,alt={图 5-1　磁场中由理想电感闭合的导轨---导体棒系统。}]{../figures/question-05.pdf}
\caption*{图 5-1　磁场中由理想电感闭合的导轨---导体棒系统。}
\end{figure}

\newpage

\section{第 6 题}\label{ux7b2c-6-ux9898}

把原恒星和现在的太阳都近似为均匀球体。原恒星早期半径远大于现在的太阳半径，可视为无穷大。已知太阳表面温度
\(T=6000\,\mathrm K\)、质量 \(M=2.0\times10^{30}\,\mathrm{kg}\)、半径
\(R=7.0\times10^8\,\mathrm m\)，万有引力常量
\(G=6.7\times10^{-11}\,\mathrm{N\,m^2\,kg^{-2}}\)，斯特藩---玻尔兹曼常量
\(\sigma=5.7\times10^{-8}\,\mathrm{W\,m^{-2}\,K^{-4}}\)。

\begin{enumerate}
\def\labelenumi{\arabic{enumi}.}
\tightlist
\item
  估算太阳由原恒星收缩到现在半径时释放的引力势能。
\item
  若按太阳当前的黑体辐射功率辐射，这些能量相当于太阳多长时间的辐射能量？
\item
  若第（1）问的能量瞬间全部转化为太阳内能，取太阳平均比热容
  \(c=10^3\,\mathrm{J\,kg^{-1}\,K^{-1}}\)，估算可使太阳达到的温度。氢聚变点火温度约为
  \(T_C=2.0\times10^7\,\mathrm K\)。假设一般恒星的平均密度和比热容与太阳相同，估算能够达到点火温度的恒星质量至少为太阳质量的百分之多少。
\end{enumerate}

\newpage

\section{第 7 题}\label{ux7b2c-7-ux9898}

\(1\,\mathrm{mol}\) 单原子分子理想气体在温度---熵图上经历圆形循环。令

\[x=\frac{T}{T_0},\qquad y=\frac{S}{S_0},\qquad S_0=2R,\]

循环满足

\[(x-2)^2+(y-2)^2=1.\]

规定气体处于 \((T_0,V_0,p_0)\) 状态时熵为 \(S_0\)。

\begin{enumerate}
\def\labelenumi{\arabic{enumi}.}
\tightlist
\item
  若循环为正循环，即气体在一个循环中对外做正功，判断图中的循环方向。
\item
  求循环效率。
\item
  求循环中的最大体积比 \(V_m/V_0\)。
\end{enumerate}

\begin{figure}[H]
\centering
\includegraphics[width=0.569\linewidth,height=\textheight,keepaspectratio,alt={图 7-1　温度---熵平面上的圆形循环。}]{../figures/question-07.pdf}
\caption*{图 7-1　温度---熵平面上的圆形循环。}
\end{figure}

\newpage

\section{第 8 题}\label{ux7b2c-8-ux9898}

假设存在磁荷为 \(g_m\) 的磁单极子，它在距自身 \(r\) 处产生的磁场为

\[\boldsymbol B=\frac{g_m\boldsymbol r}{4\pi r^3}.\]

磁单极子固定在坐标原点。电量为 \(q\) 的带电粒子从 \(z=-\infty\) 处沿直线
\(x=b\) 飞来，远处速度为
\(\boldsymbol v_0=v\hat{\boldsymbol z}\)。认为磁场造成的偏转为小量，可用未偏转的直线轨迹计算最低阶效应。

\begin{enumerate}
\def\labelenumi{\arabic{enumi}.}
\tightlist
\item
  求粒子通过磁场后的动量变化 \(\Delta\boldsymbol p\)，保留最低阶非零项。
\item
  假设偏转过程中粒子机械角动量的变化量满足玻尔式量子化条件。若由此推出的电荷最小单位应等于基本电荷
  \(e\)，求磁荷 \(g_m\)，用 \(e\) 和普朗克常量 \(h\) 表示。
\end{enumerate}

\begin{figure}[H]
\centering
\includegraphics[width=0.756\linewidth,height=\textheight,keepaspectratio,alt={图 8-1　带电粒子以碰撞参数 b 掠过磁单极子。}]{../figures/question-08.pdf}
\caption*{图 8-1　带电粒子以碰撞参数 \(b\) 掠过磁单极子。}
\end{figure}

\end{document}
